% Source: source/普通高考/2012/2012天津文.pdf, page 1, question 3.
% pdfimages -list found no embedded bitmap; the program flowchart is vector artwork.
% Grid evidence: tmp/redraw_flowchart_2012_tianjin_liberal/grid/page-1-grid100.jpg
% and tmp/redraw_flowchart_2012_tianjin_liberal/crop/q3_block_grid50.jpg.
% Node-edge list: start -> n=1,S=0 -> S=S+3^n-3^(n-1) -> n=n+1 -> test(n>=4);
% test(是) -> output S -> finish; test(否) -> return to the connector before
% the sum node. No other connectors are present.

\begin{tikzpicture}[
    >=Stealth,
    x=0.92cm,
    y=0.90cm,
    line width=0.45pt,
    line cap=round,
    line join=round,
    font=\small,
    flowbox/.style={draw, minimum height=0.68cm, inner xsep=2pt, inner ysep=1pt,
        align=center, execute at begin node={\setlength{\parindent}{0pt}}},
    terminal/.style={flowbox, rounded corners=0.14cm, minimum width=1.25cm},
    process/.style={flowbox, minimum width=2.25cm},
    decision/.style={flowbox, diamond, aspect=3.0, minimum width=3.20cm,
        minimum height=1.00cm},
    io/.style={flowbox, trapezium, minimum width=2.00cm, minimum height=0.78cm,
        trapezium left angle=70, trapezium right angle=110},
]
    \path[use as bounding box] (-4.35,-1.45) rectangle (4.45,7.55);

    \node[terminal] (start) at (0,7.10) {开始};
    \node[process, minimum width=3.45cm] (init) at (0,6.05) {$n=1,S=0$};
    \node[process, minimum width=4.00cm, minimum height=0.78cm] (sum)
        at (0,4.50) {$S=S+3^n-3^{n-1}$};
    \node[process, minimum width=2.25cm] (increment) at (0,3.15) {$n=n+1$};
    \node[decision] (test) at (0,1.75) {$n\geqslant4$};
    \node[io] (output) at (0,0.20) {输出 $S$};
    \node[terminal] (finish) at (0,-0.95) {结束};

    \draw[->] (start.south) -- (init.north);
    \draw[->] (sum.south) -- (increment.north);
    \draw[->] (increment.south) -- (test.north);
    \draw[->] (test.south) -- node[right=3pt] {是} (output.north);
    \draw[->] (output.south) -- (finish.north);

    % The vertical stem has one directed arrow into the sum node; the return
    % arrow joins that stem from the right without creating a second arrow.
    \coordinate (merge-level) at (0,5.35);
    \coordinate (merge-pre) at (0.35,5.35);
    \draw (init.south) -- (merge-level);
    \draw[->] (merge-level) -- (sum.north);

    % No branch loops right and back left to the connector before the sum.
    \draw[->] (test.east) -- node[above=2pt] {否} (3.25,1.75)
        -- (3.25,5.35) -- (merge-pre) -- (merge-level);
\end{tikzpicture}
